%杨舒云的实验报告编辑界面，使用了Huanyu Shi，2019级的模板，杨舒云在此拜谢ORZ!

%!TEX program = xelatex
\documentclass[dvipsnames, svgnames,a4paper,11pt]{article}
\input{YsySettings} 
\usepackage{lipsum}
\usepackage{adjustbox}

\begin{document}
	
	
	
	%实验报告封面	
	\begin{table}
		\renewcommand\arraystretch{1.7}
		\begin{tabularx}{\textwidth}{
				|X|X|X|X
				|X|X|X|X|}
			\hline
			\multicolumn{2}{|c|}{预习报告}&\multicolumn{2}{|c|}{实验记录}&\multicolumn{2}{|c|}{分析讨论}&\multicolumn{2}{|c|}{总成绩}\\
			\hline
			\LARGE25 & & \LARGE30 & & \LARGE25 & & \LARGE80 & \\
			\hline
		\end{tabularx}
	\end{table}
	
	\begin{table}
		\renewcommand\arraystretch{1.7}
		\begin{tabularx}{\textwidth}{|X|X|X|X|}
			\hline
			年级、专业： & 2022级 物理学 &组号： & 2-B3\\
			\hline
			姓名： & 杨舒云  & 学号： & 22344020\\
			\hline
			实验时间： & 2023/10/19 & 教师签名： & \\
			\hline
		\end{tabularx}
	\end{table}
	
	\begin{center}
		\LARGE BA4 \quad Millikan油滴实验
	\end{center}
	
	\textbf{【实验报告注意事项】}
	\begin{enumerate}
		\item 实验报告由三部分组成：
		\begin{enumerate}
			\item 预习报告：课前认真研读实验讲义，弄清实验原理；实验所需的仪器设备、用具及其使用、完成课前预习思考题；了解实验需要测量的物理量，并根据要求提前准备实验记录表格（可以参考实验报告模板，可以打印）。\textcolor{red}{\textbf{（25分）}}
			\item 实验记录：认真、客观记录实验条件、实验过程中的现象以及数据。实验记录请用珠笔或者钢笔书写并签名（\textcolor{red}{\textbf{用铅笔记录的被认为无效}}）。\textcolor{red}{\textbf{保持原始记录，包括写错删除部分，如因误记需要修改记录，必须按规范修改。}}（不得输入电脑打印，但可扫描手记后打印扫描件）；离开前请实验教师检查记录并签名。\textcolor{red}{\textbf{（30分）}}
			\item 数据处理及分析讨论：处理实验原始数据（学习仪器使用类型的实验除外），对数据的可靠性和合理性进行分析；按规范呈现数据和结果（图、表），包括数据、图表按顺序编号及其引用；分析物理现象（含回答实验思考题，写出问题思考过程，必要时按规范引用数据）；最后得出结论。\textcolor{red}{\textbf{（25分）}}
		\end{enumerate}
		\textbf{实验报告就是将预习报告、实验记录、和数据处理与分析合起来，加上本页封面。\textcolor{red}{（80分）}}
		\item 每次完成实验后的一周内交\textbf{实验报告}（特殊情况不能超过两周）。
	\end{enumerate}
	
	\textbf{【实验安全注意事项】}	
	\begin{enumerate}
		\item 实验中\textcolor{red}{\textbf{禁止用手触摸极板，注意高压危险！}}
		\item 喷油器每次喷1~2次即可，\textcolor{red}{\textbf{竖直放置、防止摔碎！}}
		\item 计时时\textcolor{red}{\textbf{眼睛要平行刻度线！}} 
	\end{enumerate}
	
	\textbf{【特别鸣谢及模板说明】}	
	
	感谢2019级学长石寰宇为本实验报告提供\LaTeX 模板。
	
	
	
	\clearpage
	\tableofcontents
	\clearpage
	
	
	
	%预习报告	
	\setcounter{section}{0}
	\section{BA4 Millikan油滴实验 \quad\heiti 预习报告}
	
	\subsection{实验目的}
	\begin{enumerate}
		\item 学习用油滴实验测量电子电荷的原理。
		\item 利用静态法和动态法观测带电油滴的运动状态， 测量不同带电油滴的电荷量， 验证电荷的不连续性， 测量电子电荷值 e。
		\item 了解 CCD 摄像机、光学系统的成像原理，了解显微测量方法以及视频信号处理技术的工程应用等。
	\end{enumerate}
	
	\subsection{仪器用具}
	\begin{table}[htbp]
		\centering
		\renewcommand\arraystretch{1.6}
		% \setlength{\tabcolsep}{10mm}
		\begin{tabular}{p{0.05\textwidth}|p{0.20\textwidth}|p{0.05\textwidth}|p{0.5\textwidth}}
			\hline
			编号& 仪器用具名称 & 数量 &  主要参数（型号，测量范围，测量精度等） \\
			\hline
			1& CCD 显微Millikan油滴仪 &1 & ZKY-MLG-6（包括视频监视器） \\
			
			2& 机油及喷雾器 &1 & -- \\
			\hline
		\end{tabular}
	\end{table}
	
	
	
	\subsection{原理概述}
	\textcolor{red}{\textbf{以下是用自己的语言描述Millikan油滴法测量电子电荷的基本原理：}}
	
	Millikan油滴实验的基本思想是通过油滴在电场和重力场下的运动来确定油滴的带电最。\cref{fig:fig1}为实验装置示意图，实验中，我们通过一个喷雾装置产生半径在亚微米级别的带电油滴，油滴在重力的作用下缓慢下落到电场区域，我们通过显微镜观察油滴在电场区域的运动情况，进而确定油滴的带电量。
	\begin{figure}[htbp]
		\centering
		\includegraphics[width=0.7\textwidth]{BA4Gra1.png}
		\caption{实验原理图（参考文献[2]）}
		\label{fig:fig1}
	\end{figure}
	\begin{enumerate}
		\item 概述：
		
		Millikan油滴法是一种用来测量电子电荷的实验方法，它是由美国物理学家罗伯特·Millikan在1909年发明的。这个方法的基本思想是，用一个喷雾器将油滴喷入一个有两块平行极板的电容器中，油滴在喷射过程中会带上一些电荷，然后在极板间受到重力和电场力的作用。通过调节极板间的电压，可以使油滴在电场中达到平衡或匀速运动的状态，从而测量出油滴的质量和带电量。由于油滴的质量和半径可以用已知的物理公式计算出来，所以只要测量出油滴的下落时间和平衡电压，就可以求出油滴所带的电荷。Millikan发现，油滴所带的电荷都是某个最小单位的整数倍，这个最小单位就是电子的电荷。
		
		\item 根据历史记载：
		\begin{itemize}
			\item Millikan的实验原理：Millikan利用两块平行极板产生一个均匀电场，然后用喷雾器向其中喷射微小的油滴。油滴在空气中带上一定的电荷，受到重力和电场力的作用。Millikan通过调节电压，使油滴在电场中达到平衡或匀速运动的状态，然后测量油滴的下落或上升的时间和速度，以及平衡时的电压。根据已知的物理量和公式，可以计算出油滴的半径、密度、质量和带电量。
			\item Millikan的实验方法：Millikan对不同大小和带电量的油滴进行了多次实验，发现油滴所带的电荷都是一个最小值的整数倍，这个最小值就是元电荷。Millikan采用了两种方法来确定元电荷：一种是直接比较不同油滴的带电量，找出它们之间的最大公约数；另一种是利用紫外线或放射源改变同一油滴的带电量，测量它们之间的差值。
			\item Millikan的实验结果：Millikan经过大量的数据分析和统计处理，得出了元电荷的平均值为 $1.592 \times 10^{-19}$ C。这个值与现代测量结果相差不大，证明了Millikan实验的精确性和可靠性。Millikan实验也为后来的原子物理学和量子力学奠定了基础，揭示了物质结构和电磁相互作用的本质。
		\end{itemize}
		
		\item 综上所述，以下是本实验的实验原理总结：
		\begin{enumerate}
			\item \textbf{电场对油滴的作用力}：首先，在实验装置中建立一个均匀的电场。这个电场会对带有电荷的油滴产生一个电场作用力，该作用力大小等于电荷乘以电场强度。
			\item \textbf{油滴的平衡和运动}：将油滴喷入电场中，观察油滴在电场中的运动。油滴会受到重力、空气阻力和电场力的影响。如果油滴的电荷是负电荷，力相抵消，使得油滴保持静止或匀速运动。
			\item \textbf{测量油滴质量和电荷}：通过观察油滴的运动和测量它的质量，可以计算出电场力和重力。
			这样，可以推导出油滴的电荷量。
			\item \textbf{电子电荷的测量}：由于油滴带有电荷，Millikan实验通过测量油滴电荷的大小，间接测定了电子的电荷量。
		\end{enumerate}		
	\end{enumerate}
	
	\textcolor{red}{\textbf{以下是预习时查找资料整理得到的更详细的实验原理：}}
	\begin{enumerate}
		\item 平衡方程
			在没有电场时，一个球形油滴在下落的过程中，会同时受到空气浮力及空气阻力。粗略地来说，速度越大，阻力越大，从而最终比将达到受力平衡。由粘性流体的知识可知，在速度较低时（严格而言应该是雷诺数较低时），球形物体收到的阻力正比于其运动速度 $v$，比例系数为 $6\pi R \eta$，亦即:
			\begin{equation*}
				F_\mathrm{d} = -6 \pi r v \eta
			\end{equation*}
			这里 $r$ 是油滴半径，$v$ 是油滴速度，$\eta$ 为流体的动力粘度，其大小反映了粘性流体的流速场存在一定梯度时能量耗散的大小，而负号表示其总是与运动速度方向相反。上式是流体力学中一个著名的公式，名为 Stokes 公式。
			
			将浮力和重力项加入到平衡方程中，可得:
			\begin{equation*}
				m_1 g - m_2 g - 6 \pi r v \eta = 0
			\end{equation*}
			这里 $m_1$ 为油滴质量，$m_2$ 为同体积空气的质量（浮力项），$v_\mathrm {f}$ 为下落时的平衡速度（这里用下标 f 以与后面上升平衡速度 $v_\mathrm {r}$ 相区分）。
			
			考虑到油滴体积较小，此时空气不再能直接视为连续流体，因此 Stokes 公式中的动力粘度需要进行修正。修正的形式为:
			\begin{equation*}
				F_\mathrm{d} = -6 \pi r v \eta \left(1 + \frac{b}{pr}\right)
			\end{equation*}
			这里 $b$ 为一常系数，这里取 $b = 8.23 \times 10^ {-3} \mathrm {N} / \mathrm {m}$。$p$ 为大气压，$r$ 为液滴半径。若令:
			\begin{equation*}
				K = \left(1 + \frac{b}{pr}\right)
			\end{equation*}
			则平衡方程可以写成:
			\begin{equation*}
				m_1 g - m_2 g - K6 \pi r v_\mathrm {f} \eta = 0
			\end{equation*}
			
			\begin{figure}[htbp]
				\centering
				\includegraphics[width=0.7\textwidth]{BA4Gra2.png}
				\caption{油滴在电场中受力示意图（参考文献[3]）}
				\label{fig:fig2}
			\end{figure}
			
			另一方面，若假设油滴带正电，在施加一定的强度的向上的电场时，油滴可能会处于静止平衡的状态。此时的平衡方程为:
			\begin{equation*}
				q E_\mathrm {s} + m_2 g - m_1 g - K6 \pi r v_\mathrm {s} \eta = 0
			\end{equation*}
			这里 $E_\mathrm {s}$ 表示油滴静止时的平衡电场，$q$ 为油滴带电量。
			
			若向上的电场继续增加，油滴将会以匀速上升的方式达到平衡。设平衡速度为 $v_\mathrm {r}$，对应电压为 $E_\mathrm {r}$，则平衡方程为:
			\begin{equation*}
				q E_\mathrm {r} + m_2 g - m_1 g - K6 \pi r v_\mathrm {r} \eta = 0
			\end{equation*}
			
			实际上，本实验中的核心方程就是下落平衡方程，静止平衡方程以及上升平衡方程。可知，其中主要的未知量只是油滴电量 $q$ 以及油滴半径 $r$。由于这里有三个方程，因此只需取两个就可以解出所有未知。后面的推导中将看到，通过下落平衡方程的近似方程可以得到油滴半径的粗估值，之后分别带入静止平衡方程以及上升平衡方程即可得到油滴电量，前者称为静态法，后者称为动态法。
			
			\item 油滴半径的估计
			做粗略计算时 (可以看成将 Stokes 定律修正中的修正项取零阶近似), 可以直接由下落平衡方程解出油滴半径。 另一方面, 由于油滴质量不容易测量, 实验上习惯采用油滴密度来计算质量, 即:
			\begin{align*}
				m_1 &= \dfrac{4}{3}\pi r^3\rho_1 \\
				m_2 &= \dfrac{4}{3}\pi r^3\rho_2
			\end{align*}
			将油滴质量及空气质量带入下落平衡方程后, 可解出油滴半径的估计值 (这里用 $r_0$ 代表估计值):
			$$r\approx r_0=\left[\dfrac{9v_f\eta}{(\rho_1-\rho_2)g}\right]^{1/2}$$
			
			\item 油滴电量的测量
			考虑到实验中采用均匀电场, 设施加电压为 U, 板间距为 d, 则电场强度可表示为:
			$$E=\dfrac{U}{d}$$
			另一方面, 实验中油滴的下落距离 s 是固定的, 因此可以用下落时间来计算下落速度:
			$$v=\dfrac{s}{t}$$
			除电场强度, 下落速度外, 将油滴半径, 阻尼系数, 油滴质量, 空气质量也带入上升平衡方程, 可得:
			$$q=9\sqrt{2}\pi d\left[\dfrac{(\eta s)^3}{(\rho_1-\rho_2)g}\right]^{1/2}\left(\dfrac{1}{1+\frac{b}{pr_0}}\right)^{3/2}\cdot\dfrac{1}{U_r}\left(\dfrac{1}{t_f}+\dfrac{1}{t_r}\right)\left(\dfrac{1}{t_f}\right)^{1/2}$$
			这里 $U_r$ 表示上升平衡时的电压. 另一方面, 虽然也可以直接从静态平衡方程推出静态电压, 但方便的做法是将静态平衡视为平衡速度为零的动态平衡, 亦即 $t_r \to \infty$。 将该式带入到动态法的电荷方程中, 即可得到静态法的电压:
			$$q=9\sqrt{2}\pi d\left[\dfrac{(\eta s)^3}{(\rho_1-\rho_2)g}\right]^{1/2}\left(\dfrac{1}{1+\frac{b}{pr_0}}\right)^{3/2}\cdot\dfrac{1}{U_s}\left(\dfrac{1}{t_f}\right)^{3/2}$$
			这里 $U_s$ 表示静止平衡时的电压。		
	\end{enumerate}
	
	除以上之外，还有一些历史记载可以给我们这个实验的启示：
	\begin{itemize}
		\item 1981年，阿兰·富兰克林（Allan Franklin）研究了Millikan的实验记录本，发现Millikan在记录本中对其观察结果进行打分，从“一般”到“最好”。根据记录本，Millikan在1913年发表的论文依据的是140次观察，然而他把其中49次观察的数据舍弃不用，只根据91次他认为较好的观察结果的数据进行计算。但是，在论文中，Millikan却声称该论文“代表了所有的油滴实验”。如果Millikan把所有的观察数据都包括进去，虽然不会影响其结果，却会加大误差。
		
		\item 我记得当时用的办法是假定已知一个电子带电量为$1.6\times 10^{-19}$，再用油滴带电量除电子电量看有几个电子（此处经常除不断，我记得我还除出来一个$17.5$，极其离谱）然后就得出这个油滴带$18$个电子，然后用$\frac{18-17.5}{18}\times 100\%=3$
		得出误差是$3\%$，实验成功。虽然计算上最后获得了成功，但是回头看这数据处理实在是觉得良心痛。
		
		\item Millikan油滴实验60年后，史学家发现，Millikan一共向外公布了58次观测数据，而他本人一共做过140次观测。他在实验中通过预先估测，去掉了那些他认为有偏差，误差大的数据。也就是说，如果你抽奖10000次，100次中奖，9900次没中，然后去掉9900个错误答案，剩下100个中奖率100。
		油滴实验到现在依然是重要的物理实验，无数莘莘学子前仆后继地做，甚至导师会善意提醒数据有偏差没关系，不要编数。
		
		\item 58组符合理论的数据证明了理论和实验的正确，剩下的82组数据，都是别人的错。Millikan做油滴实验的结果与今天的测量值相比有一些偏差，原因是Millikan使用了错误的空气粘滞系数。但是人们一开始没有发现，因为做实验的人发现数据比Millikan的实验值大一些，便怀疑自己做错了，一通改动之后再重新来一组新数据，直到“与Millikan看齐”为止。这一切是整个科学的耻辱。密令油滴实验反而是英雄，向人们展示了科学的严谨，以及科学工作者不可告人的潜规则。
		
		\item 查理·费曼曾经在1974年，于加州理工学院的一场毕业典礼演说中叙述了“草包族科学”（Cargo cult science）：从过往的经验，我们学到了如何应付一些自我欺骗的情况。举个例子，Millikan做了个油滴实验，量出了电子的带电量，得到一个今天我们知道是不大对的答案。他的资料有点偏差，因为他用了个不准确的空气粘滞系数数值。于是，如果你把在Millikan之后、进行测量电子带电量所得到的资料整理一下，就会发现一些很有趣的现象：把这些资料跟时间画成坐标图，你会发现这个人得到的数值比Millikan的数值大一点点，下一个人得到的资料又再大一点点，下一个又再大上一点点，最后，到了一个更大的数值才稳定下来。为什么他们没有在一开始就发现新数值应该较高？——这件事令许多相关的科学家惭愧脸红——因为显然很多人的做事方式是：当他们获得一个比Millikan数值更高的结果时，他们以为一定哪里出了错，他们会拼命寻找，并且找到了实验有错误的原因。另一方面，当他们获得的结果跟Millikan的相仿时，便不会那么用心去检讨。因此，他们排除了所谓相差太大的数据，不予考虑。我们现在已经很清楚那些伎俩了，因此再也不会犯同样的毛病。
	\end{itemize}
	
	
	\subsection{实验前思考题}
	\begin{question}
		Millikan利用油滴测定电子电荷的基本原理是什么？
	\end{question}
	按要求，此题不作答，详见\textbf{原理概述}部分。
	
	\begin{question}
		什么是静态(平衡)测量法和动态(非平衡)测量法？两种方法有何不同与优缺点？测量中需注意哪些问题？
	\end{question}
	按要求，此题不作答，详见讲义。
	
	\begin{question}
		为什么必须保证油滴在测量范围内做匀速运动或静止？怎样控制油滴运动？
	\end{question}
	按实验原理，要求油滴受力平衡，即油滴做匀速运动或者处于静止状态。
	
	控制油滴运动：首先调节电压为零，喷出油滴后观察仪器的显示器，寻找一个下落速度适中的油滴，下落时间介于 20 到 30 秒为宜，后调节平衡电压使该油滴处于平衡状态，则油滴符合要求。
	
	具体来说，有以下几种想法。
	
	调整电场强度：通过调整电场强度，可以控制电场对油滴的作用力，使其与重力达到平衡，油滴保持匀速运动或静止。
	
	调整油滴的电荷量：通过控制油滴的电荷量，可以影响电场对油滴的作用力，进而调整油滴的运动状态。增加或减少油滴的电荷量可以使其达到匀速运动或静止状态。
	
	观察和微调：实验者需要仔细观察油滴的运动状态，并根据观察结果微调电场强度或油滴的电荷量，以确保油滴处于所需的运动状态。
	
	%实验记录	
	\clearpage
	\begin{table}
		\renewcommand\arraystretch{1.7}
		\centering
		\begin{tabularx}{\textwidth}{|X|X|X|X|}
			\hline
			专业： & 物理学 & 年级： & 2022级 \\
			\hline
			姓名： & 杨舒云 & 学号： & 22344020\\
			\hline
			室温： & 26摄氏度 & 实验地点： & 实验室509 \\
			\hline
			学生签名：& 杨舒云 & 评分： &\\
			\hline
			实验时间：& 2023/10/19 & 教师签名：&\\
			\hline
		\end{tabularx}
	\end{table}
	
	\section{BA4 Millikan油滴实验 \quad\heiti 实验记录}
	\subsection{实验内容、步骤与结果}
	\subsubsection{仪器检查和测量准备}
	\begin{enumerate}
		\item 调整油滴实验仪：小心取下 ZKY-MLG-6-CCD 显微Millikan油滴仪的金属盖放置一边。水平调整。实验仪使用。CCD 成像系统调整。
		\item 学习控制油滴在视场中的运动：选择合适的油滴；平衡电压的确认；控制油滴的运动；熟悉各按键的操作后，可以尝试测量数据 
	\end{enumerate}
	
	\subsubsection{平衡法}
	\begin{enumerate}
		\item 将“0V/工作”按键状态切换至“工作”，红色指示灯点亮；将“平衡/提升”按键置于“平衡”。将平衡电压调整为 200到400V 之间任意值。
		\item 通过喷雾口向油滴盒内喷入油雾，此时监视器上将出现大量运动的油滴。选取合适的油滴，仔细调整平衡电压 U，使其平衡在最上面的格线上。
		\item 将“0V/工作”状态按键切换至“0V”，此时油滴开始下落，当油滴下落到有“0”标记的格线时，立即按下计时开始键，计时器启动，开始记录油滴下落时间 t。
		\item 当油滴下落至有距离标记的格线时（本实验：1.6），立即按下计时结束键，计时器停止计时，油滴立即静止，“0V/工作”按键自动切换至“工作”。通过“确认”按键将这次测量的“平衡电压 U 和匀速下落时间 t”记录在监视器屏幕上。
		\item 将“平衡/提升”按键置于“提升”，油滴将向上运动，当回到最上面的格线时，将“平衡/提升”键切换至平衡，油滴停止上升，重新调整平衡电压。（注意：如果此处的平衡电压发生了突变，则该油滴得到或失去了电子。这次测量不能采用，需要重新寻找油滴进行测量。）
		\item 当对同一个油滴的 5 次测量完成后，按“确认”键，实验仪内部的单片机将根据参数自动计算出 5 次测量的平均平衡电压和平均匀速下落时间，并根据参数自动计算并显示出该油滴的电荷量 q。
		\item 记录实验界面显示的数值（包括每次用的平衡电压，下落时间， 5 次平均电压，5 次平均时间，测量电荷量的平均值）。
		\item 重复上述步骤，要求测量 5 颗不同的油滴，测量每颗油滴的电荷量 q。
		
		此处只罗列将在数据分析当中用于不确定评定的数据，分别来自我自己的油滴4、戴鹏辉22344016的油滴1和朱政鑫22344019的油滴1、
		如\cref{tab:tab1}所示。
		\begin{table}[h]
			\centering
			\caption{部分数据展示}
			\label{tab:tab1}
			\resizebox{\textwidth}{!}{%
			\begin{tabular}{|c|c|c|c|c|c|c|c|c|}
				\hline
				Y油滴4 & 平衡电压$U_g/V$ & 下落时间$t_g/s$ & D油滴1 & 平衡电压$U_g/V$ & 下落时间$t_g/s$ & Z油滴1 & 平衡电压$U_g/V$ & 下落时间$t_g/s$\\
				\hline
				& 266 & 6.00 & & 222 & 15.99 & & 253 & 10.01 \\
				& 266 & 6.11 & & 223 & 16.18 & & 260 & 10.27 \\
				& 266 & 6.09 & & 223 & 16.23 & & 260 & 10.00 \\
				& 266 & 5.96 & & 223 & 16.40 & & 260 & 10.00 \\
				& 266 & 5.99 & & 223 & 16.01 & & 257 & 10.15 \\
				\hline
			\end{tabular}
		}
		\end{table}
		
		事实上，仔细观察可以发现我在原始数据中的记录是出现了笔误的，为了证明我没有伪造数据，
		我将拍摄的实验结果展示出来（\cref{fig:figA3}）。
		\begin{figure}[htbp]
			\centering
			\includegraphics[width=0.9\textwidth]{BA4GraA3.jpg}
			\caption{原始数据3}
			\label{fig:figA3}
		\end{figure}
		 
		
		其它测量数据见数据展示部分，附：原始数据如\cref{fig:figA1}所示。
		\begin{figure}[htbp]
			\centering
			\includegraphics[width=0.9\textwidth]{BA4GraA1.jpg}
			\caption{原始数据1}
			\label{fig:figA1}
		\end{figure}
		
		数据分析见\textbf{实验数据分析}部分。
	\end{enumerate}
	
	\clearpage
	\subsubsection{动态法}
	与平衡法不同，动态法分两步完成。第一步是油滴下落过程，其操作同平衡法。完成第一步后，如果对本次测量结果满意，则可以按下确认键保存这个步骤的测量结果，如果不满意，则可以删除（删除方法见前面所述）。
	
	第一步骤完成后，油滴处于距离标志格线以下（本实验：1.6）。通过“0V/工作”键、“平衡/提升”键配合，使油滴下偏距离“1.6”标志格线一定距离。调节“电压调节”旋钮加大电压，使油滴上升，当油滴到达“1.6”标志格线时，立即按下计时开始键，计时器开始计时；当油滴上升到“0”标记格线时，再次按下计时键，停止计时，但油滴继续上升，再次调节“电压调节”旋钮使油滴平衡于“0”格线以上（图 6），按下“确认”键保存本次实验结果。
	
	重复以上步骤完成 5 次完整实验，按下确认键，实验仪自动计算出实验结果。动态测量法是分别测出下落时间、提升时间及提升电压，并代入式（16）计算求得油滴带电量 q。
			
	测量数据见数据展示部分，附：原始数据如\cref{fig:figA1}、\cref{fig:figA2}所示。
	\begin{figure}[htbp]
		\centering
		\includegraphics[width=0.9\textwidth]{BA4GraA2.jpg}
		\caption{原始数据2}
		\label{fig:figA2}
	\end{figure}
	
	数据分析见\textbf{实验数据分析}部分。
	
	\subsubsection{数据展示}
	由于实验的结果过于糟糕，我和同学（戴鹏辉22344016）将截止到第八周的所有实验数据全部整理了出来，得到了如\cref{tab:tab2}所示的实验测量结果（这个表是能通过放大看清的）。
	\begin{table}[h]
		\centering
		\caption{测量结果}
		\label{tab:tab2}
		\resizebox{\textwidth}{!}{%
			\begin{tabular}{cccc|cccc|cccc|cccc|cccc|cccc}
				\hline
				测量编号 & 油滴带电量/$10^{-19}C$ & 电子个数 & 元电荷量/$10^{-19}C$ & 测量编号 & 油滴带电量/$10^{-19}C$ & 电子个数 & 元电荷量/$10^{-19}C$ & 测量编号 & 油滴带电量/$10^{-19}C$ & 电子个数 & 元电荷量/$10^{-19}C$ & 测量编号 & 油滴带电量/$10^{-19}C$ & 电子个数 & 元电荷量/$10^{-19}C$ & 测量编号 & 油滴带电量/$10^{-19}C$ & 电子个数 & 元电荷量/$10^{-19}C$ & 测量编号 & 油滴带电量/$10^{-19}C$ & 电子个数 & 元电荷量/$10^{-19}C$ \\
				\hline
				1 & 11.107 & 7 & 1.587 & 65 & 1.510 & 1 & 1.510 & 129 & 11.140 & 7 & 1.591 & 193 & 10.720 & 7 & 1.531 & 257 & 11.865 & 7 & 1.695 & 321 & 1.551 & 1 & 1.551 \\
				2 & 1.505 & 1 & 1.505 & 66 & 32.188 & 20 & 1.609 & 130 & 4.569 & 3 & 1.523 & 194 & 4.757 & 3 & 1.586 & 258 & 11.577 & 7 & 1.654 & 322 & 1.529 & 1 & 1.529 \\
				3 & 15.614 & 10 & 1.561 & 67 & 17.687 & 11 & 1.608 & 131 & 3.047 & 2 & 1.524 & 195 & 12.116 & 8 & 1.515 & 259 & 18.132 & 11 & 1.648 & 323 & 9.376 & 6 & 1.563 \\
				4 & 3.019 & 2 & 1.510 & 68 & 1.445 & 1 & 1.445 & 132 & 1.547 & 1 & 1.547 & 196 & 13.999 & 9 & 1.555 & 260 & 13.407 & 8 & 1.676 & 324 & 6.337 & 4 & 1.584 \\
				5 & 3.075 & 2 & 1.538 & 69 & 14.757 & 9 & 1.640 & 133 & 3.022 & 2 & 1.511 & 197 & 12.088 & 8 & 1.511 & 261 & 13.863 & 9 & 1.540 & 325 & 2.922 & 2 & 1.461 \\
				6 & 34.913 & 22 & 1.587 & 70 & 18.968 & 12 & 1.581 & 134 & 3.038 & 2 & 1.519 & 198 & 1.700 & 1 & 1.700 & 262 & 8.589 & 5 & 1.718 & 326 & 20.650 & 13 & 1.588 \\
				7 & 3.540 & 2 & 1.770 & 71 & 6.788 & 4 & 1.697 & 135 & 1.503 & 1 & 1.503 & 199 & 4.767 & 3 & 1.589 & 263 & 66.023 & 41 & 1.610 & 327 & 1.528 & 1 & 1.528 \\
				8 & 3.680 & 2 & 1.840 & 72 & 1.729 & 1 & 1.729 & 136 & 8.291 & 5 & 1.658 & 200 & 1.605 & 1 & 1.605 & 264 & 2.745 & 2 & 1.373 & 328 & 13.893 & 9 & 1.544 \\
				9 & 2.990 & 2 & 1.495 & 73 & 5.097 & 3 & 1.699 & 137 & 9.606 & 6 & 1.601 & 201 & 15.617 & 10 & 1.562 & 265 & 14.607 & 9 & 1.623 & 329 & 1.543 & 1 & 1.543 \\
				10 & 4.560 & 3 & 1.520 & 74 & 6.368 & 4 & 1.592 & 138 & 1.614 & 1 & 1.614 & 202 & 29.769 & 19 & 1.567 & 266 & 6.340 & 4 & 1.585 & 330 & 6.328 & 4 & 1.582 \\
				11 & 4.320 & 3 & 1.440 & 75 & 14.817 & 9 & 1.646 & 139 & 3.107 & 2 & 1.554 & 203 & 15.603 & 10 & 1.560 & 267 & 1.568 & 1 & 1.568 & 331 & 4.528 & 3 & 1.509 \\
				12 & 0.857 & 1 & 0.857 & 76 & 13.812 & 9 & 1.535 & 140 & 3.159 & 2 & 1.580 & 204 & 3.016 & 2 & 1.508 & 268 & 3.020 & 2 & 1.510 & 332 & 1.498 & 1 & 1.498 \\
				13 & 0.917 & 1 & 0.917 & 77 & 3.125 & 2 & 1.563 & 141 & 12.600 & 8 & 1.575 & 205 & 1.615 & 1 & 1.615 & 269 & 22.300 & 14 & 1.593 & 333 & 1.842 & 1 & 1.842 \\
				14 & 1.432 & 1 & 1.432 & 78 & 1.618 & 1 & 1.618 & 142 & 3.052 & 2 & 1.526 & 206 & 2.086 & 1 & 2.086 & 270 & 5.877 & 4 & 1.469 & 334 & 6.583 & 4 & 1.646 \\
				15 & 1.142 & 1 & 1.142 & 79 & 4.559 & 3 & 1.520 & 143 & 1.513 & 1 & 1.513 & 207 & 20.002 & 13 & 1.539 & 271 & 4.298 & 3 & 1.433 & 335 & 6.461 & 4 & 1.615 \\
				16 & 6.099 & 4 & 1.525 & 80 & 5.064 & 3 & 1.688 & 144 & 3.112 & 2 & 1.556 & 208 & 12.114 & 8 & 1.514 & 272 & 5.369 & 3 & 1.790 & 336 & 1.683 & 1 & 1.683 \\
				17 & 9.290 & 6 & 1.548 & 81 & 7.749 & 5 & 1.550 & 145 & 1.578 & 1 & 1.578 & 209 & 10.762 & 7 & 1.537 & 273 & 19.514 & 12 & 1.626 & 337 & 5.998 & 4 & 1.500 \\
				18 & 2.988 & 2 & 1.494 & 82 & 17.473 & 11 & 1.588 & 146 & 1.563 & 1 & 1.563 & 210 & 1.568 & 1 & 1.568 & 274 & 2.592 & 2 & 1.296 & 338 & 1.559 & 1 & 1.559 \\
				19 & 4.835 & 3 & 1.612 & 83 & 1.731 & 1 & 1.731 & 147 & 6.153 & 4 & 1.538 & 211 & 7.920 & 5 & 1.584 & 275 & 7.863 & 5 & 1.573 & 339 & 5.826 & 4 & 1.457 \\
				20 & 9.253 & 6 & 1.542 & 84 & 1.715 & 1 & 1.715 & 148 & 1.565 & 1 & 1.565 & 212 & 7.912 & 5 & 1.582 & 276 & 9.089 & 6 & 1.515 & 340 & 1.430 & 1 & 1.430 \\
				21 & 10.655 & 7 & 1.522 & 85 & 3.392 & 2 & 1.696 & 149 & 4.680 & 3 & 1.560 & 213 & 7.543 & 5 & 1.509 & 277 & 1.558 & 1 & 1.558 & 341 & 3.080 & 2 & 1.540 \\
				22 & 4.641 & 3 & 1.547 & 86 & 1.692 & 1 & 1.692 & 150 & 10.991 & 7 & 1.570 & 214 & 16.595 & 10 & 1.660 & 278 & 3.069 & 2 & 1.535 & 342 & 6.669 & 4 & 1.667 \\
				23 & 6.162 & 4 & 1.541 & 87 & 1.727 & 1 & 1.727 & 151 & 4.583 & 3 & 1.528 & 215 & 9.150 & 6 & 1.525 & 279 & 7.782 & 5 & 1.556 & 343 & 11.587 & 7 & 1.655 \\
				24 & 13.462 & 8 & 1.683 & 88 & 7.571 & 5 & 1.514 & 152 & 4.686 & 3 & 1.562 & 216 & 3.020 & 2 & 1.510 & 280 & 6.048 & 4 & 1.512 & 344 & 10.046 & 6 & 1.674 \\
				25 & 3.054 & 2 & 1.527 & 89 & 4.489 & 3 & 1.496 & 153 & 3.079 & 2 & 1.540 & 217 & 1.516 & 1 & 1.516 & 281 & 6.001 & 4 & 1.500 & 345 & 10.322 & 6 & 1.720 \\
				26 & 16.476 & 10 & 1.648 & 90 & 10.950 & 7 & 1.564 & 154 & 3.215 & 2 & 1.608 & 218 & 6.213 & 4 & 1.553 & 282 & 2.983 & 2 & 1.492 & 346 & 7.180 & 4 & 1.795 \\
				27 & 40.700 & 25 & 1.628 & 91 & 6.063 & 4 & 1.516 & 155 & 3.040 & 2 & 1.520 & 219 & 1.559 & 1 & 1.559 & 283 & 3.002 & 2 & 1.501 & 347 & 15.796 & 10 & 1.580 \\
				\hline
			\end{tabular}
		}
		
		\resizebox{\textwidth}{!}{%
			\begin{tabular}{cccc|cccc|cccc|cccc|cccc|cccc}
				\hline
				测量编号 & 油滴带电量/$10^{-19}C$ & 电子个数 & 元电荷量/$10^{-19}C$ & 测量编号 & 油滴带电量/$10^{-19}C$ & 电子个数 & 元电荷量/$10^{-19}C$ & 测量编号 & 油滴带电量/$10^{-19}C$ & 电子个数 & 元电荷量/$10^{-19}C$ & 测量编号 & 油滴带电量/$10^{-19}C$ & 电子个数 & 元电荷量/$10^{-19}C$ & 测量编号 & 油滴带电量/$10^{-19}C$ & 电子个数 & 元电荷量/$10^{-19}C$ & 测量编号 & 油滴带电量/$10^{-19}C$ & 电子个数 & 元电荷量/$10^{-19}C$ \\
				\hline
				28 & 9.286 & 6 & 1.548 & 92 & 12.195 & 8 & 1.524 & 156 & 14.523 & 9 & 1.614 & 220 & 4.805 & 3 & 1.602 & 284 & 4.300 & 3 & 1.433 & 348 & 2.902 & 2 & 1.451 \\
				29 & 6.030 & 4 & 1.508 & 93 & 7.527 & 5 & 1.505 & 157 & 7.885 & 5 & 1.577 & 221 & 3.075 & 2 & 1.538 & 285 & 3.019 & 2 & 1.510 & 349 & 6.033 & 4 & 1.508 \\
				30 & 7.491 & 5 & 1.498 & 94 & 6.243 & 4 & 1.561 & 158 & 3.107 & 2 & 1.554 & 222 & 4.828 & 3 & 1.609 & 286 & 1.493 & 1 & 1.493 & 350 & 29.539 & 18 & 1.641 \\
				31 & 4.592 & 3 & 1.531 & 95 & 6.258 & 4 & 1.565 & 159 & 3.160 & 2 & 1.580 & 223 & 7.583 & 5 & 1.517 & 287 & 1.477 & 1 & 1.477 & 351 & 24.009 & 15 & 1.601 \\
				32 & 4.570 & 3 & 1.523 & 96 & 6.166 & 4 & 1.542 & 160 & 3.085 & 2 & 1.543 & 224 & 8.267 & 5 & 1.653 & 288 & 4.572 & 3 & 1.524 & 352 & 2.970 & 2 & 1.485 \\
				33 & 7.406 & 5 & 1.481 & 97 & 6.087 & 4 & 1.522 & 161 & 10.730 & 7 & 1.533 & 225 & 1.243 & 1 & 1.243 & 289 & 4.537 & 3 & 1.512 & 353 & 1.452 & 1 & 1.452 \\
				34 & 3.470 & 2 & 1.735 & 98 & 6.124 & 4 & 1.531 & 162 & 1.600 & 1 & 1.600 & 226 & 3.882 & 2 & 1.941 & 290 & 4.644 & 3 & 1.548 & 354 & 1.469 & 1 & 1.469 \\
				35 & 3.637 & 2 & 1.819 & 99 & 3.242 & 2 & 1.621 & 163 & 4.586 & 3 & 1.529 & 227 & 14.419 & 9 & 1.602 & 291 & 1.533 & 1 & 1.533 & 355 & 7.446 & 5 & 1.489 \\
				36 & 6.248 & 4 & 1.562 & 100 & 3.993 & 2 & 1.997 & 164 & 4.471 & 3 & 1.490 & 228 & 15.449 & 10 & 1.545 & 292 & 4.115 & 3 & 1.372 & 356 & 1.488 & 1 & 1.488 \\
				37 & 4.543 & 3 & 1.514 & 101 & 10.750 & 7 & 1.536 & 165 & 13.614 & 9 & 1.513 & 229 & 3.168 & 2 & 1.584 & 293 & 6.458 & 4 & 1.615 & 357 & 9.317 & 6 & 1.553 \\
				38 & 4.746 & 3 & 1.582 & 102 & 4.917 & 3 & 1.639 & 166 & 83.100 & 52 & 1.598 & 230 & 10.781 & 7 & 1.540 & 294 & 9.144 & 6 & 1.524 & 358 & 11.003 & 7 & 1.572 \\
				39 & 9.275 & 6 & 1.546 & 103 & 7.669 & 5 & 1.534 & 167 & 6.411 & 4 & 1.603 & 231 & 6.155 & 4 & 1.539 & 295 & 7.589 & 5 & 1.518 & 359 & 3.081 & 2 & 1.541 \\
				40 & 44.177 & 28 & 1.578 & 104 & 2.406 & 2 & 1.203 & 168 & 23.330 & 15 & 1.555 & 232 & 10.858 & 7 & 1.551 & 296 & 1.488 & 1 & 1.488 & 360 & 7.766 & 5 & 1.553 \\
				41 & 16.925 & 11 & 1.539 & 105 & 12.280 & 8 & 1.535 & 169 & 15.207 & 10 & 1.521 & 233 & 7.681 & 5 & 1.536 & 297 & 6.701 & 4 & 1.675 & 361 & 3.112 & 2 & 1.556 \\
				42 & 1.390 & 1 & 1.390 & 106 & 3.050 & 2 & 1.525 & 170 & 17.306 & 11 & 1.573 & 234 & 3.090 & 2 & 1.545 & 298 & 1.534 & 1 & 1.534 & 362 & 9.450 & 6 & 1.575 \\
				43 & 3.896 & 2 & 1.948 & 107 & 6.197 & 4 & 1.549 & 171 & 66.736 & 42 & 1.589 & 235 & 3.060 & 2 & 1.530 & 299 & 6.600 & 4 & 1.650 & 363 & 9.277 & 6 & 1.546 \\
				44 & 2.006 & 1 & 2.006 & 108 & 3.050 & 2 & 1.525 & 172 & 17.471 & 11 & 1.588 & 236 & 32.499 & 20 & 1.625 & 300 & 3.486 & 2 & 1.743 & 364 & 4.689 & 3 & 1.563 \\
				45 & 3.167 & 2 & 1.584 & 109 & 6.023 & 4 & 1.506 & 173 & 6.284 & 4 & 1.571 & 237 & 12.734 & 8 & 1.592 & 301 & 6.164 & 4 & 1.541 & 365 & 4.676 & 3 & 1.559 \\
				46 & 7.379 & 5 & 1.476 & 110 & 10.823 & 7 & 1.546 & 174 & 24.004 & 15 & 1.600 & 238 & 9.411 & 6 & 1.569 & 302 & 9.736 & 6 & 1.623 & 366 & 4.771 & 3 & 1.590 \\
				47 & 9.678 & 6 & 1.613 & 111 & 3.109 & 2 & 1.555 & 175 & 2.411 & 2 & 1.205 & 239 & 19.858 & 12 & 1.655 & 303 & 1.493 & 1 & 1.493 & 367 & 3.500 & 2 & 1.750 \\
				48 & 7.805 & 5 & 1.561 & 112 & 3.036 & 2 & 1.518 & 176 & 3.530 & 2 & 1.765 & 240 & 19.800 & 12 & 1.650 & 304 & 3.110 & 2 & 1.555 & 368 & 11.970 & 7 & 1.710 \\
				49 & 7.912 & 5 & 1.582 & 113 & 3.054 & 2 & 1.527 & 177 & 4.834 & 3 & 1.611 & 241 & 12.140 & 8 & 1.518 & 305 & 10.710 & 7 & 1.530 & 369 & 2.001 & 1 & 2.001 \\
				50 & 1.538 & 1 & 1.538 & 114 & 13.394 & 8 & 1.674 & 178 & 1.486 & 1 & 1.486 & 242 & 7.717 & 5 & 1.543 & 306 & 11.122 & 7 & 1.588 & 370 & 4.582 & 3 & 1.527 \\
				51 & 2.996 & 2 & 1.498 & 115 & 7.260 & 4 & 1.815 & 179 & 3.194 & 2 & 1.597 & 243 & 3.002 & 2 & 1.501 & 307 & 4.201 & 3 & 1.400 & 371 & 4.652 & 3 & 1.551 \\
				52 & 8.682 & 5 & 1.736 & 116 & 1.561 & 1 & 1.561 & 180 & 9.225 & 6 & 1.538 & 244 & 1.561 & 1 & 1.561 & 308 & 4.148 & 3 & 1.383 & 372 & 1.562 & 1 & 1.562 \\
				53 & 3.500 & 2 & 1.750 & 117 & 7.757 & 5 & 1.551 & 181 & 3.093 & 2 & 1.547 & 245 & 1.640 & 1 & 1.640 & 309 & 2.985 & 2 & 1.493 & 373 & 1.565 & 1 & 1.565 \\
				\hline
			\end{tabular}
		}
		
		\resizebox{\textwidth}{!}{%
			\begin{tabular}{cccc|cccc|cccc|cccc|cccc|cccc}
				\hline
				测量编号 & 油滴带电量/$10^{-19}C$ & 电子个数 & 元电荷量/$10^{-19}C$ & 测量编号 & 油滴带电量/$10^{-19}C$ & 电子个数 & 元电荷量/$10^{-19}C$ & 测量编号 & 油滴带电量/$10^{-19}C$ & 电子个数 & 元电荷量/$10^{-19}C$ & 测量编号 & 油滴带电量/$10^{-19}C$ & 电子个数 & 元电荷量/$10^{-19}C$ & 测量编号 & 油滴带电量/$10^{-19}C$ & 电子个数 & 元电荷量/$10^{-19}C$ & 测量编号 & 油滴带电量/$10^{-19}C$ & 电子个数 & 元电荷量/$10^{-19}C$ \\
				\hline
				54 & 2.998 & 2 & 1.499 & 118 & 3.066 & 2 & 1.533 & 182 & 3.124 & 2 & 1.562 & 246 & 1.531 & 1 & 1.531 & 310 & 5.348 & 3 & 1.783 & 374 & 3.029 & 2 & 1.514 \\
				55 & 9.482 & 6 & 1.580 & 119 & 6.347 & 4 & 1.587 & 183 & 4.703 & 3 & 1.568 & 247 & 6.489 & 4 & 1.622 & 311 & 1.580 & 1 & 1.580 & 375 & 3.085 & 2 & 1.542 \\
				56 & 4.792 & 3 & 1.597 & 120 & 3.117 & 2 & 1.559 & 184 & 1.468 & 1 & 1.468 & 248 & 1.552 & 1 & 1.552 & 312 & 6.178 & 4 & 1.545 & 376 & 2.993 & 2 & 1.497 \\
				57 & 1.482 & 1 & 1.482 & 121 & 4.503 & 3 & 1.501 & 185 & 1.535 & 1 & 1.535 & 249 & 5.556 & 3 & 1.852 & 313 & 3.170 & 2 & 1.585 & 377 & 11.540 & 7 & 1.648 \\
				58 & 1.488 & 1 & 1.488 & 122 & 11.103 & 7 & 1.586 & 186 & 6.212 & 4 & 1.553 & 250 & 1.526 & 1 & 1.526 & 314 & 1.508 & 1 & 1.508 & 378 & 3.004 & 2 & 1.502 \\
				59 & 8.640 & 5 & 1.728 & 123 & 3.049 & 2 & 1.525 & 187 & 3.166 & 2 & 1.583 & 251 & 5.866 & 3 & 1.955 & 315 & 6.401 & 4 & 1.600 & 379 & 5.591 & 3 & 1.864 \\
				60 & 8.565 & 5 & 1.713 & 124 & 3.116 & 2 & 1.558 & 188 & 5.623 & 3 & 1.874 & 252 & 5.652 & 3 & 1.884 & 316 & 13.799 & 9 & 1.533 & 380 & 4.511 & 3 & 1.504 \\
				61 & 6.477 & 4 & 1.619 & 125 & 10.931 & 7 & 1.561 & 189 & 6.448 & 4 & 1.612 & 253 & 1.577 & 1 & 1.577 & 317 & 14.263 & 9 & 1.585 & 381 & 1.531 & 1 & 1.531 \\
				62 & 9.244 & 6 & 1.541 & 126 & 10.974 & 7 & 1.568 & 190 & 6.122 & 4 & 1.531 & 254 & 9.530 & 6 & 1.588 & 318 & 6.344 & 4 & 1.586 & 382 & 7.907 & 5 & 1.581 \\
				63 & 3.596 & 2 & 1.798 & 127 & 1.555 & 1 & 1.555 & 191 & 4.507 & 3 & 1.502 & 255 & 4.605 & 3 & 1.535 & 319 & 1.583 & 1 & 1.583 & 383 & 5.519 & 3 & 1.840 \\
				64 & 6.028 & 4 & 1.507 & 128 & 8.267 & 5 & 1.653 & 192 & 3.090 & 2 & 1.545 & 256 & 1.520 & 1 & 1.520 & 320 & 4.682 & 3 & 1.561 & 384 & 4.544 & 3 & 1.515 \\
				\hline
			\end{tabular}
		}
		\end{table}
	
	\clearpage
	\subsection{原始数据记录}
	实验记录本上的原始数据见\cref{fig:figA1}、\cref{fig:figA2}。
	
	实验台桌面整理见\textbf{附件}部分(\cref{fig:figA4})。
	
	\subsection{实验过程中遇到的问题记录}
	\begin{enumerate}
		\item 我个人感觉的主要问题是无法确定油滴是处于平衡状态，即油滴的运动难以控制；
		\item 除此之外，油滴的选取也很困难；
		\item 计时时会因为我的反应时间带来较大的误差；
		\item 实验结果与理论值的偏差相当大。
	\end{enumerate}
	
	
	
	%分析与讨论	
	\clearpage
	\begin{table}
		\renewcommand\arraystretch{1.7}
		\begin{tabularx}{\textwidth}{|X|X|X|X|}
			\hline
			专业：& 物理学 &年级：& 2022级\\
			\hline
			姓名： & 杨舒云 & 学号：& 22344020\\
			\hline
			日期：& 2023/10/19 & 评分： &\\
			\hline
		\end{tabularx}
	\end{table}
	
	\section{BA4 Millikan油滴实验 \quad\heiti 分析与讨论}
	\subsection{实验数据分析}
	\textcolor{red}{\textbf{特别说明：由于本实验报告全部计算内容由计算机软件（Excel、Mathematica和Python）完成，因此可能会出现数据输入错误、精度丢失、运算结果出错等问题。若老师批改时发现自己算的和我算的不一样，请先检查自己的计算是否有问题；若还是对不上，请不要怀疑本实验报告的原创性，请联系我获得Excel原始表格及本实验报告的\LaTeX 代码，以便进一步确定错误来源。}}
	\subsubsection{按讲义上的方法处理实验数据}
	\begin{enumerate}
		\item 按Millikan的方法对数据进行处理
		
		初步分析\cref{tab:tab2}得到的数据呈现出以下结果：	
		\begin{center}
			\begin{tabular}{|c|c|c|}
				\hline
				元电荷电量平均值 & 单次测量标准差 & 平均值的标准差\\
				1.564594367 & 0.113043457 & 0.000294384\\
				\hline
				数据总数 & 大于1.6 & 小于1.6\\
				384 & 101 & 282 \\
				\hline
			\end{tabular}
		\end{center}
		
		按照Millikan的思路，我将上述数据进行了分级（优良差），选取其中等级为优的数据，如\cref{tab:tab3}所示。
		\begin{table}[h]
			\centering
			\caption{优选结果}
			\label{tab:tab3}
			\resizebox{0.9\textwidth}{!}{%
			\begin{tabular}{|c|c|c|c|c|c|c|c|}
				\hline
				测量编号 & 油滴带电量/$10^{-19}C$ & 电子个数 & 元电荷量/$10^{-19}C$ & 测量编号 & 油滴带电量/$10^{-19}C$ & 电子个数 & 元电荷量/$10^{-19}C$ \\
				\hline
				1 & 11.107 & 7 & 1.587 & 32 & 66.023 & 41 & 1.610 \\
				2 & 34.913 & 22 & 1.587 & 33 & 14.607 & 9 & 1.623 \\
				3 & 4.835 & 3 & 1.612 & 34 & 6.340 & 4 & 1.585 \\
				4 & 40.700 & 25 & 1.628 & 35 & 22.300 & 14 & 1.593 \\
				5 & 4.746 & 3 & 1.582 & 36 & 19.514 & 12 & 1.626 \\
				6 & 3.167 & 2 & 1.584 & 37 & 6.458 & 4 & 1.615 \\
				7 & 9.678 & 6 & 1.613 & 38 & 9.736 & 6 & 1.623 \\
				8 & 3.189 & 2 & 1.595 & 39 & 4.885 & 3 & 1.628 \\
				9 & 1.584 & 1 & 1.584 & 40 & 9.606 & 6 & 1.601 \\
				10 & 4.746 & 3 & 1.582 & 41 & 1.614 & 1 & 1.614 \\
				11 & 6.324 & 4 & 1.581 & 42 & 3.215 & 2 & 1.608 \\
				12 & 4.757 & 3 & 1.586 & 43 & 14.523 & 9 & 1.614 \\
				13 & 4.767 & 3 & 1.589 & 44 & 3.160 & 2 & 1.580 \\
				14 & 1.605 & 1 & 1.605 & 45 & 1.600 & 1 & 1.600 \\
				15 & 1.615 & 1 & 1.615 & 46 & 83.100 & 52 & 1.598 \\
				16 & 7.920 & 5 & 1.584 & 47 & 6.411 & 4 & 1.603 \\
				17 & 7.912 & 5 & 1.582 & 48 & 66.736 & 42 & 1.589 \\
				18 & 4.805 & 3 & 1.602 & 49 & 17.471 & 11 & 1.588 \\
				19 & 4.828 & 3 & 1.609 & 50 & 24.004 & 15 & 1.600 \\
				20 & 14.419 & 9 & 1.602 & 51 & 21.032 & 13 & 1.618 \\
				21 & 3.168 & 2 & 1.584 & 52 & 48.470 & 30 & 1.616 \\
				22 & 32.499 & 20 & 1.625 & 53 & 1.620 & 1 & 1.620 \\
				23 & 12.734 & 8 & 1.592 & 54 & 1.598 & 1 & 1.598 \\
				24 & 9.630 & 6 & 1.605 & 55 & 6.337 & 4 & 1.584 \\
				25 & 20.802 & 13 & 1.600 & 56 & 20.650 & 13 & 1.588 \\
				26 & 32.188 & 20 & 1.609 & 57 & 6.461 & 4 & 1.615 \\
				27 & 17.687 & 11 & 1.608 & 58 & 24.009 & 15 & 1.601 \\
				28 & 18.968 & 12 & 1.581 & 59 & 4.771 & 3 & 1.590 \\
				29 & 1.618 & 1 & 1.618 & 60 & 11.201 & 7 & 1.600 \\
				30 & 17.473 & 11 & 1.588 & 61 & 21.001 & 13 & 1.615 \\
				31 & 3.242 & 2 & 1.621 &  &  &  &  \\
				\hline
			\end{tabular}
		}
		\end{table}
				
		\item 计算法
		
		实验测量 5 颗油滴，记录每颗油滴的电荷量 $q_i$,再用电荷 $e$ 的理论值去除，
		对商四舍五入取整后得到每颗油滴所带电子个数 $n_i$，之后再求 $q_i/n_i=q_e^{(i)}$得到每次测量
		的基本电荷，再求出 5 次测量的均值$\bar{q_e}$ ，与理论值进行对比。
		
		根据\cref{tab:tab3}，我们可计算得到如下结果：
		
		均值：$\bar{q_e}=1.601C$；
		
		标准差：$\sigma=0.014C$；
		
		实验标准差：$\sigma_{\bar{q_e}}=0.002C$；
		
		相对误差：$\eta=0.90\%$。
				
		\item 作图法
		
		得到 $q_i$ 和对应的 $n_i$ 后，以 $q_i$ 为纵坐标， $n_i$ 为横坐标作图，拟合得到的直线
		斜率即为基本电荷 e 的测量值，与理论值进行对比。如\cref{fig:fig3}所示。
		\begin{figure}[htbp]
			\centering
			\includegraphics[width=0.7\textwidth]{BA4Gra3.png}
			\caption{作图法求元电荷量}
			\label{fig:fig3}
		\end{figure}
		
		由此求出斜率即元电荷量：$q_e=(1.60303 ± 0.00218)C$。
		
		\item 由电量分布图说明油滴带电量的不连续性
		
		下面我将给出油滴的带电量分布图。此图的经典作法是先将电量划分成若干小区间，然后统计落在每个区间内的油滴数量，最后作出柱形图。这种方式作图不仅要求有非常大量的数据，并且作出来的图连续性极差，不太直观。因此，在作图时可以把每一个数据点看成一个正态分布，所有数据点对应的正态分布叠加后得到的便是我们需要求的油滴带电量分布。
		\begin{figure}[htbp]
			\centering
			\includegraphics[width=0.7\textwidth]{BA4Gra4.png}
			\caption{油滴带电量分布图}
			\label{fig:fig4}
		\end{figure}
		
		如\cref{fig:fig4}所示，可以很明显地看到数个峰，这很好地说明了油滴带电量确实是不连续的。
		
		最后补充一点关于正态分布的取法：对于任一个数据，服从$N( \mu, \sigma^2)$，$\mu$取其值，$\sigma$取统一的不确定度（详见\textbf{误差分析}部分）。
	\end{enumerate}
	
	\clearpage
	\subsubsection{误差分析}
	\begin{enumerate}
		\item 计算相对误差
		
		按上一节的数据处理，已经得到计算法的结果
		
		标准差：$\sigma=0.014C$；
		
		实验标准差：$\sigma_{\bar{q_e}}=0.002C$；
		
		相对误差：$\eta=0.90\%$。
		
		因此不对此再赘述了，下面将选取\cref{tab:tab1}的数据进行具体的不确定度评定。
		
		\item 关于误差来源的讨论
		
		事实上，误差来源在实验过程中我已经有了相当好的体验。首先一大误差来源便是计时：单片机的计时误差基本可以忽略，但测量者和机器的按钮都有一定的反应延时，有些此测量者还可能会尝试通过提前按按钮来补偿机器的反应延时。这一部分误差对于同一个测量者的多次测量而言通常是稳定并且均值非零的，因此需要计入系统误差。考志到测量者在自己能够反应的时间尺度上是会对机器的延时进行补偿的，因此计时的不确定度实际上就是测量者的反应时间，这里取这个反应时间为 0.3s。其次，我们很难保证油滴是真正平衡了的，这是另一部分误差，不过我个人认为它可以归结到随机误差当中。
		
		按照对整个384组数据的分析，最后的结果小于1.6的足足有282组，这说明了有目前尚未明白的系统误差存在。
		
		除以上之外，我个人认为还应该考虑电压测量的系统误差，但由于相关参数的缺乏，暂不考虑这一部分误差。
		\item 不确定度评定
		
		按上述误差分析，作如下不确定度评定：
		
		主要作A类评定取：$u_A$为实验标准差；B类评定主要考虑计时因素，故取$\sqrt{3}u_B=0.3s$。
		
		由此计算合成不确定度及展伸不确定度。
		
		\item 不确定度报告
		
		\begin{equation*}
			q_e=\bar{q_e}\pm U_{q_e}\approx( 1.60 \pm  0.02 )\times10^{-19}C
		\end{equation*} 
		其中展伸不确定度$U_{q_e}={k}\times{u_{q_e}}\approx 0.02\times10^{-19} C$ 是由合成标准不确定度$u_{q_e}\approx 0.01\times10^{-19} C$ 和$k\approx4.26$确定的，k是依据置信概率$P=99.73\%$和自由度$\nu\approx8.09$查t分布表得到的。
								
	\end{enumerate}
	
	\subsubsection{尝试一种新的数据处理方式}
	在我看来，本实验对数据的处理方法“实验测量 5 颗油滴，记录每颗油滴的电荷量 $q_i$,再用电荷 $e$ 的理论值去除，
	对商四舍五入取整后得到每颗油滴所带电子个数 $n_i$”科学性不足，正如在原理概述中所说的“虽然计算上最后获得了成功，但是回头看这数据处理实在是觉得良心痛”。因此，我尝试着使用一些别的处理数据的方式（虽然我觉得这个新方式依旧缺乏充足的科学性）。
	\begin{enumerate}
		\item 具体方法
		
		对于大量的电荷量的实验数据（假设已经从小到大排列），分别计算相邻的两个数据的差值得到,观察一下发现它们要么接近于0，要么很像是某一个数的整数倍，于是乎我们猜测电荷有一个最小单位记作$e$；先给出一个试探的解$e_1$，让它等于一系列数据中非零（接近于0的当作0）的最小数,现在我们可以肯定的是，元电荷的真实值不会大于试探解$e_1$，并且当$e_j>2e_1$时，我们也可认为$e_j$是不可能的，但对于其它试探解都是有可能的解，现在得到了一系列试探解${e_n}$；对于某一个试探解$e_k$，我们用原始数据除以它得到，四舍五入之后得到$n_k$，根据$n_k$的结果，我们对各${e_n}$进行加权平均从而得到最终结果。
		
		\item 实际操作
		
		根据上述方法，我编写了如附件中展示的第二段代码来实现。并得到了一系列试解（单位均为$10^{-19}C$）：
		$\Delta e=0.77, e_1=1.357, e_2=1.886, e_3= 2.287$。
		
		接下来进一步处理数据：$\bar{e}=0.5\times1.357+0.4\times1.886+0.1\times2.287=1.662$，可以看出，实际上的结果并不好。（在笔者看来，可以类比着附件中第一段代码的方法用正态分布对上述三个试解进行拟合）。
		
		若由$\Delta e=0.77$剔除掉$e_3= 2.287$，则得到$\bar{e}=1.622$，这是一个较好的结果。
	\end{enumerate}
	
	\subsubsection{结论梳理}
	\begin{enumerate}
		\item \textbf{元电荷量大概在$1.6\times10^{-19}C$左右；}
		\item \textbf{油滴带电量是不连续的；}
		\item \textbf{实验中存在着未知的系统误差，目前已确定的主要误差来源是时间和电压的测量。}
	\end{enumerate}
	
	\clearpage
	\subsection{实验后思考题}
	\begin{question}
		试分析本实验的主要误差来源。
	\end{question}
	实验误差来源我已经在\textbf{误差分析}部分有了详细的讨论，现罗列于下：事实上，误差来源在实验过程中我已经有了相当好的体验。首先一大误差来源便是计时：单片机的计时误差基本可以忽略，但测量者和机器的按钮都有一定的反应延时，有些此测量者还可能会尝试通过提前按按钮来补偿机器的反应延时。这一部分误差对于同一个测量者的多次测量而言通常是稳定并且均值非零的，因此需要计入系统误差。考志到测量者在自己能够反应的时间尺度上是会对机器的延时进行补偿的，因此计时的不确定度实际上就是测量者的反应时间，这里取这个反应时间为 0.3s。其次，我们很难保证油滴是真正平衡了的，这是另一部分误差，不过我个人认为它可以归结到随机误差当中。事实上，按照对整个384组数据的分析，最后的结果小于1.6的足足有282组，这说明了有目前尚未明白的系统误差存在。除以上之外，我个人认为还应该考虑电压测量的系统误差，但由于相关参数的缺乏，暂不考虑这一部分误差。
	
	除以上之外，这里再补充一部分误差来源：
	\begin{itemize}
		\item 空气和油的密度：空气密度会随温度、湿度和压力而变化，油的密度可能不太准确。 这些因素会影响油滴的半径和质量的计算，从而影响电荷的计算。
		
		\item 电场强度：在两个平行板之间，电场强度并不完全恒定或等于电压除以板间距。由于板间距和板面积相当，所以边缘效应可能会改变电场强度的分布，导致实际的电场强度与测量值有偏差。
		
		\item 光学系统：光学系统的质量会影响油滴的可见性和清晰度。如果油滴不够亮或者不够清晰，那么测量它们运动的时间和距离就会有误差. 此外，显微镜系统的校准也会影响距离的测量。
		
		\item Cunningham修正因子：Cunningham修正因子是一个考虑空气分子对油滴运动影响的因子。如果忽略这个因子，那么油滴的粘滞阻力就会被低估，从而导致电荷的计算有偏差。
	\end{itemize}
		
	\begin{question}
		Millikan油滴实验的总结（油滴筛选、跟踪、测量等）（经验分享；体会；感想；讨论；建议等）
	\end{question}
	\begin{enumerate}
		\item \textbf{实验概述}：Millikan油滴实验的目的是测量电子的电荷量。它利用油滴在电场中的运动来确定电子的电荷。
		
		\item \textbf{油滴筛选}：选择合适的油滴非常重要。一般使用云母片将油滴悬浮在空气中，然后使用一个细雾喷雾器将油滴带入实验室。在筛选油滴时，确保选择那些大小均匀、稳定的油滴，以减小误差。
		
		\item \textbf{跟踪油滴}：追踪油滴的运动是实验的关键。通常使用显微镜来观察油滴在电场中的运动。这需要细心和耐心，因为油滴的轨迹可能会变化。确保记录油滴的位置和速度，并进行多次观察以减小误差。
		
		\item \textbf{测量电场强度}：测量电场强度是实验的另一个关键步骤。使用库仑定律和已知电压来计算电场强度，然后使用这些数据来确定电子电荷量。
		
		\item \textbf{数据分析}：实验数据的分析非常重要。使用已知的物理原理，例如库仑定律和牛顿力学，来分析数据并计算电子电荷量。
		
		\item \textbf{经验分享和感想}：Millikan油滴实验需要精确性和细致的观察。在实验中，我们真正体会到了物理学实验的重要性，以及如何应用理论知识来解决实际问题。这个实验也向我们展示了科学方法的力量，以及如何逐步改进实验设计和减小误差。
		
		\item \textbf{讨论}：Millikan油滴实验的结果对于原子结构理论的发展产生了深远的影响。它为科学家提供了测量电子电荷量的关键数据，有助于建立了现代原子模型。此外，实验还引发了对微观世界的更深入研究，如量子力学的兴起。
		
		\item \textbf{建议}：在进行Millikan油滴实验时，要保持仔细、耐心和精确。进行多次实验以获取更多数据，以减小误差。此外，与同学合作，讨论实验结果和数据分析，可以提高实验的可靠性。
	\end{enumerate}
	
	
	
	\clearpage
	\section{BA4 Millikan油滴实验 \quad\heiti 参考文献与附件}
	\subsection{参考文献}
	[1]沈韩.基础物理实验.——北京：科学出版社，2015.2 ISBN：978-7-03-043311-4 
	
	[2]Wikipedia contributors. (2023, April 6). Oil drop experiment. In Wikipedia, The Free Encyclopedia. Retrieved 12:18, October 18, 2023, from [https://en.wikipedia.org/w/index.php?title=Oil drop experiment\&oldid=1016715200].
	
	[3]Leonard, E. (2018, January 22). Millikan’s Oil Drop Experiment: A Simple Explanation. Owlcation. Retrieved 12:18, October 18, 2023, from [https://owlcation.com/stem/Millikans-Oil-Drop-Experiment].  
	
	\subsection{附件}
	\begin{figure}[htbp]
		\centering
		\includegraphics[width=0.9\textwidth]{BA4GraA4.jpg}
		\caption{实验台桌面整理}
		\label{fig:figA4}
	\end{figure}
	
	\begin{lstlisting}[style=pythonstyle,caption=代码记录1]
		import numpy as np
		import pandas as pd
		import matplotlib.pyplot as plt
		
		num = pd.read_excel('D:\JupyterNotebooks\BA4Data.xlsx')
		num_list=num['油滴带电量/10^-19C'].tolist()
		
		for num in num_list:
		if num > 30:
		num_list.remove(num)
		num_array = np.array(num_list)
		#print(num_list)
		
		# 每个数据点的均值和标准差，可以根据需要进行调整
		mean_values = num_array
		std_deviation = 0.05*1.602
		
		# 创建x轴的数值范围
		x = np.linspace(0, 100, 1000)
		
		# 初始化一个空的概率密度曲线
		pdf_curve = np.zeros(x.shape)
		
		# 将每个正态分布叠加到概率密度曲线上
		for mean in mean_values:
		pdf_curve += (1 / (std_deviation * np.sqrt(2 * np.pi))) * np.exp(-((x - mean) ** 2) / (2 * std_deviation ** 2))
		
		# 绘制概率密度曲线
		plt.xlim(0, 30)
		plt.plot(x, pdf_curve, label="Oil Droplet Charge Distribution Curve")
		plt.xlabel("Oil Droplet Charge")
		plt.ylabel("Probability Density")
		plt.legend()
		plt.show()
	\end{lstlisting}
	
	\begin{lstlisting}[style=pythonstyle,caption=代码记录1]
		import pandas as pd
		
		data = pd.read_excel('D:\JupyterNotebooks\BA4Data.xlsx')
		#print(data)
		data_list = data['油滴带电量/10^-19C'].tolist()
		#print(data_list)
		
		# 定义一个函数，输入为一系列数据，输出为任意两者差值的绝对值的最小值
		def min_diff(data):
		# 如果数据为空或只有一个元素，返回None
		if not data or len(data) == 1:
		return None
		# 对数据进行排序
		data.sort()
		# 初始化最小差值为第一个元素和第二个元素的差值的绝对值
		min_diff = 100
		# 遍历数据中的相邻元素，更新最小差值
		for i in range(1, len(data) - 1):
		diff = abs(data[i] - data[i + 1])
		if ((diff > 2.87) and (diff < min_diff)):
		min_diff = diff
		# 返回最小差值
		return min_diff
		
		print(min_diff(data_list)) 
		
	\end{lstlisting}

\end{document}